%%%
%%% Radical "⾛"
%%%
\section*{Radical 156: ``⾛''}\addcontentsline{toc}{section}{Radical 156: ⾛}\addcontentsline{loh}{figure}{\#\#\#\# 156: ⾛}

%%%%%%%%%% 走 %%%%%%%%%%
\subsection*{走}\addcontentsline{loh}{figure}{走}

\begin{Entry}{走}{7}{⾛}[Kangxi 156]
  \begin{Phonetics}{走}{zou3}[][HSK 1]
    \definition{v.}{andar; caminhar | correr | mover; movimentar; deslocar | sair; partir; ir embora | visitar; fazer uma visita; (entre amigos e familiares) troca de visitas | passar por; atravessar; ultrapassar | vazar; revelar; divulgar | afastar"-se do original; alterar ou perder a forma, o sabor, a cor, etc. originais}
  \synonymref{离}{li2}
  \synonymref{跑}{pao3}
  \synonymref{行}{xing2}
  \end{Phonetics}
\end{Entry}

\begin{Entry}{走开}{7,4}{⾛,⼶}
  \begin{Phonetics}{走开}{zou3kai1}[][HSK 2]
    \definition{v.}{ir embora; fugir; ir para outro lugar}
  \synonymref{离开}{li2/kai1}
  \synonymref{远离}{yuan3li2}
  \antonymref{到来}{dao4lai2}
  \antonymref{归来}{gui1lai2}
  \antonymref{回来}{hui2lai5}
  \antonymref{靠拢}{kao4long3}
  \end{Phonetics}
\end{Entry}

\begin{Entry}{走去}{7,5}{⾛,⼛}
  \begin{Phonetics}{走去}{zou3 qu4}
    \definition{v.}{caminhar até (para)}
  \end{Phonetics}
\end{Entry}

\begin{Entry}{走电}{7,5}{⾛,⽥}
  \begin{Phonetics}{走电}{zou3dian4}
    \definition{s.}{Dialeto: fuga de eletricidade; fuga elétrica}
  \synonymref{漏电}{lou4dian4}
  \synonymref{跑电}{pao3dian4}
  \end{Phonetics}
\end{Entry}

\begin{Entry}{走后门}{7,6,3}{⾛,⼝,⾨}
  \begin{Phonetics}{走后门}{zou3hou4men2}[][HSK 7-9]
    \definition{v.}{puxar os cordões; entrar pela porta dos fundos, é uma metáfora para o uso de meios ilícitos, como favores ou subornos, para atingir um determinado objetivo por meio de conexões internas}
  \end{Phonetics}
\end{Entry}

\begin{Entry}{走过}{7,6}{⾛,⾡}
  \begin{Phonetics}{走过}{zou3guo4}[][HSK 2]
    \definition{v.}{passar por; perambular}
  \end{Phonetics}
\end{Entry}

\begin{Entry}{走过场}{7,6,6}{⾛,⾡,⼟}
  \begin{Phonetics}{走过场}{zou3guo4chang3}[][HSK 7-9]
    \definition{adj.}{Coloquial: como uma mera formalidade; superficial; sem rigor; sem impacto}
    \definition{v.}{Coloquial: cumprir formalidades; fazer as coisas mecanicamente}
  \end{Phonetics}
\end{Entry}

\begin{Entry}{走投无路}{7,7,4,13}{⾛,⼿,⽆,⾜}
  \begin{Phonetics}{走投无路}{zou3tou2-wu2lu4}[][HSK 7-9]
    \definition{expr.}{ser encurralado; estar no limite; estar em apuros; estar desesperado; estar sem saída; estar em uma situação desesperadora; estar em um dilema sem saída; estar em um impasse; estar em uma situação difícil; estar preso em um dilema; entre a cruz e a espada; não encontrar saída; chegar a um beco sem saída; chegar ao limite; levar alguém ao limite; sentir"-se encurralado; estar em uma situação sem saída; não ter a quem recorrer; não ter saída; não ter como pedir ajuda; entrar em um beco sem saída; não saber para que lado ir; pobre e completamente indefeso}
  \synonymref{一筹莫展}{yi4chou2-mo4zhan3}
  \antonymref{鹏程万里}{peng2cheng2-wan4li3}
  \antonymref{千方百计}{qian1fang1-bai3ji4}
  \antonymref{如愿以偿}{ru2yuan4yi3chang2}
  \end{Phonetics}
\end{Entry}

\begin{Entry}{走秀}{7,7}{⾛,⽲}
  \begin{Phonetics}{走秀}{zou3xiu4}
    \definition{s.}{desfile de moda}
    \definition{v.}{andar na passarela (em um desfile de moda)}
  \end{Phonetics}
\end{Entry}

\begin{Entry}{走私}{7,7}{⾛,⽲}
  \begin{Phonetics}{走私}{zou3/si1}[][HSK 6]
    \definition{v.+compl.}{contrabandear; ter um caso ilícito; violar regulamentos alfandegários; evadir inspeções alfandegárias; transportar mercadorias ilegalmente para dentro e para fora do país}
  \end{Phonetics}
\end{Entry}

\begin{Entry}{走近}{7,7}{⾛,⾡}
  \begin{Phonetics}{走近}{zou3jin4}[][HSK 7-9]
    \definition{v.}{aproximar"-se; aproximar"-se de}
  \end{Phonetics}
\end{Entry}

\begin{Entry}{走进}{7,7}{⾛,⾡}
  \begin{Phonetics}{走进}{zou3jin4}[][HSK 2]
    \definition{v.}{entrar}
  \end{Phonetics}
\end{Entry}

\begin{Entry}{走势}{7,8}{⾛,⼒}
  \begin{Phonetics}{走势}{zou3shi4}
    \definition{s.}{tendência | direção; alinhamento}
  \end{Phonetics}
\end{Entry}

\begin{Entry}{走卒}{7,8}{⾛,⼗}
  \begin{Phonetics}{走卒}{zou3zu2}
    \definition{s.}{peão; capacho; lacaio; fantoche | lacaio (de malfeitor) | peão (ou seja, soldado de infantaria) | servo}
  \antonymref{主人}{zhu3ren2}
  \antonymref{主宰}{zhu3zai3}
  \end{Phonetics}
\end{Entry}

\begin{Entry}{走弯路}{7,9,13}{⾛,⼸,⾜}
  \begin{Phonetics}{走弯路}{zou3wan1lu4}[][HSK 7-9]
    \definition{v.}{fazer um desvio; seguir um percurso em ziguezague; pegar um caminho errado | tomar uma rota indireta | perder tempo usando um método inadequado}
  \end{Phonetics}
\end{Entry}

\begin{Entry}{走鬼}{7,9}{⾛,⿁}
  \begin{Phonetics}{走鬼}{zou3 gui3}
    \definition{s.}{vendedor ambulante sem licença; camelô}
  \end{Phonetics}
\end{Entry}

\begin{Entry}{走索}{7,10}{⾛,⽷}
  \begin{Phonetics}{走索}{zou3suo3}
    \definition{v.}{Acrobacia: andar na corda bamba; dança da corda; caminhada na corda}
  \seealsoref{走绳}{zou3sheng2}
  \end{Phonetics}
\end{Entry}

\begin{Entry}{走廊}{7,11}{⾛,⼴}
  \begin{Phonetics}{走廊}{zou3lang2}[][HSK 7-9]
    \definition[条]{s.}{corredor; passagem; uma passagem acima do solo sob os beirais; uma passagem coberta entre edifícios | corredor; passagem (uma faixa de terra que liga duas áreas maiores); uma metáfora para uma estreita faixa de terra que liga duas grandes regiões}
  \end{Phonetics}
\end{Entry}

\begin{Entry}{走绳}{7,11}{⾛,⽷}
  \begin{Phonetics}{走绳}{zou3sheng2}
    \definition{v.}{andar na corda bamba; dança da corda; caminhada na corda; um tipo de acrobacia em que atores andam para frente e para trás em uma corda suspensa e realizam diversas ações}
  \seealsoref{走索}{zou3suo3}
  \end{Phonetics}
\end{Entry}

\begin{Entry}{走路}{7,13}{⾛,⾜}
  \begin{Phonetics}{走路}{zou3lu4}[][HSK 1]
    \definition{v.}{caminhar; ir a pé; andar em pé sobre a terra | sair; ir embora; partir}
  \synonymref{步行}{bu4xing2}
  \synonymref{撤离}{che4li2}
  \synonymref{赶路}{gan3lu4}
  \synonymref{行走}{xing2zou3}
  \synonymref{走开}{zou3kai1}
  \antonymref{奔跑}{ben1pao3}
  \antonymref{静止}{jing4zhi3}
  \antonymref{留下}{liu2xia4}
  \antonymref{停留}{ting2liu2}
  \antonymref{站立}{zhan4li4}
  \end{Phonetics}
\end{Entry}

%%%%%%%%%% 赴 %%%%%%%%%%
\subsection*{赴}\addcontentsline{loh}{figure}{赴}

\begin{Entry}{赴}{9}{⾛}
  \begin{Phonetics}{赴}{fu4}[][HSK 7-9]
    \definition{v.}{ir para | comparecer}
  \end{Phonetics}
\end{Entry}

%%%%%%%%%% 赶 %%%%%%%%%%
\subsection*{赶}\addcontentsline{loh}{figure}{赶}

\begin{Entry}{赶}{10}{⾛}
  \begin{Phonetics}{赶}{gan3}[][HSK 3]
    \definition*{s.}{Sobrenome: Gan}
    \definition{prep.}{por; até; até que; até quando; introduzir o momento em que algo aconteceu, indicando que se espera até um determinado momento}
    \definition{v.}{ultrapassar; alcançar | perseguir; correr atrás; tentar alcançar; dar uma corrida; acelerar ou intensificar  | dirigir; conduzir | expulsar; afugentar; afastar | encontrar; deparar"-se com; esbarrar em; acontecer; encontrar"-se em (uma situação); aproveitar"-se de (uma oportunidade) | ir para; participar (atividades com horário marcado)}
  \end{Phonetics}
\end{Entry}

\begin{Entry}{赶上}{10,3}{⾛,⼀}
  \begin{Phonetics}{赶上}{gan3shang4}[][HSK 6]
    \definition{v.}{alcançar; manter o ritmo com; acompanhar alguém ou o padrão do planejador | chegar a tempo para; ter tempo suficiente; não ser tarde demais | encontrar; topar com; cruzar com; encontrar"-se com; acontecer de encontrar; encontrar algo, em um determinado momento ou oportunidade}
  \end{Phonetics}
\end{Entry}

\begin{Entry}{赶不上}{10,4,3}{⾛,⼀,⼀}
  \begin{Phonetics}{赶不上}{gan3bu5shang4}[][HSK 6]
    \definition{v.}{ficar para trás; ser incapaz de alcançar; não conseguir alcançar; não conseguir acompanhar | ser tarde demais (para fazer algo); (não) existir tempo suficiente (para fazer algo) |  deixar de ter; ser incapaz de encontrar ou ter a chance de encontrar; não encontrar; não encontrar (boa oportunidade) | não poder ser comparado a}
  \end{Phonetics}
\end{Entry}

\begin{Entry}{赶忙}{10,6}{⾛,⼼}
  \begin{Phonetics}{赶忙}{gan3mang2}[][HSK 6]
    \definition{adv.}{imediatamente; com pressa; às pressas; rapidamente}
  \end{Phonetics}
\end{Entry}

\begin{Entry}{赶早}{10,6}{⾛,⽇}
  \begin{Phonetics}{赶早}{gan3zao3}
    \definition{adv.}{o mais cedo possível; antes que seja tarde demais; o mais breve possível; na primeira oportunidade; quanto antes, melhor}
  \seealsoref{赶早儿}{gan3zao3r5}
  \synonymref{趁早}{chen4zao3}
  \synonymref{及早}{ji2zao3}
  \synonymref{提早}{ti2zao3}
  \end{Phonetics}
\end{Entry}

\begin{Entry}{赶早儿}{10,6,2}{⾛,⽇,⼉}
  \begin{Phonetics}{赶早儿}{gan3zao3r5}
    \definition{adv.}{o mais cedo possível; antes que seja tarde demais}
  \end{Phonetics}
\end{Entry}

\begin{Entry}{赶快}{10,7}{⾛,⼼}
  \begin{Phonetics}{赶快}{gan3kuai4}[][HSK 3]
    \definition{adv.}{rapidamente; imediatamente; aproveite o momento e acelere o ritmo}
  \seealsoref{赶紧}{gan3jin3}
  \seealsoref{赶忙}{gan3mang2}
  \seealsoref{尽快}{jin3kuai4}
  \synonymref{赶紧}{gan3jin3}
  \synonymref{赶忙}{gan3mang2}
  \synonymref{急忙}{ji2mang2}
  \synonymref{尽快}{jin3kuai4}
  \synonymref{快捷}{kuai4jie2}
  \synonymref{快速}{kuai4su4}
  \synonymref{连忙}{lian2mang2}
  \synonymref{马上}{ma3shang4}
  \synonymref{迅速}{xun4su4}
  \antonymref{怠慢}{dai4man4}
  \antonymref{缓慢}{huan3man4}
  \antonymref{拖延}{tuo1yan2}
  \end{Phonetics}
\end{Entry}

\begin{Entry}{赶走}{10,7}{⾛,⾛}
  \begin{Phonetics}{赶走}{gan3zou3}
    \definition{v.}{ir embora; sair de carro | expulsar; mandar embora; mandar alguém embora; expulsar alguém pela porta; mostrar a alguém a porta; expulsar alguém pela porta; expulsar alguém de casa; chutar alguém escada abaixo; colocar alguém contra a porta; expulsar alguém a pontapés; ver as costas de alguém; perseguir; expulsar alguém}
  \synonymref{驱逐}{qu1zhu2}
  \antonymref{接纳}{jie1na4}
  \antonymref{收留}{shou1liu2}
  \end{Phonetics}
\end{Entry}

\begin{Entry}{赶到}{10,8}{⾛,⼑}
  \begin{Phonetics}{赶到}{gan3dao4}[][HSK 3]
    \definition{v.}{correr (para algum lugar); apressar"-se}
  \synonymref{到达}{dao4da2}
  \synonymref{抵达}{di3da2}
  \antonymref{出发}{chu1fa1}
  \antonymref{离开}{li2/kai1}
  \end{Phonetics}
\end{Entry}

\begin{Entry}{赶往}{10,8}{⾛,⼻}
  \begin{Phonetics}{赶往}{gan3wang3}[][HSK 7-9]
    \definition{v.}{apressar"-se para (algum lugar)}
  \end{Phonetics}
\end{Entry}

\begin{Entry}{赶赴}{10,9}{⾛,⾛}
  \begin{Phonetics}{赶赴}{gan3fu4}[][HSK 7-9]
    \definition{v.}{apressar"-se para; correr para | apressar"-se}
  \end{Phonetics}
\end{Entry}

\begin{Entry}{赶紧}{10,10}{⾛,⽷}
  \begin{Phonetics}{赶紧}{gan3jin3}[][HSK 3]
    \definition{adv.}{apressadamente; precipitadamente; às pressas; significa agir imediatamente, sem demora}
  \seealsoref{赶快}{gan3kuai4}
  \seealsoref{连忙}{lian2mang2}
  \seealsoref{一下子}{yi2xia4zi5}
  \synonymref{赶快}{gan3kuai4}
  \synonymref{赶忙}{gan3mang2}
  \synonymref{急忙}{ji2mang2}
  \synonymref{尽快}{jin3kuai4}
  \synonymref{快速}{kuai4su4}
  \synonymref{连忙}{lian2mang2}
  \synonymref{马上}{ma3shang4}
  \antonymref{从容}{cong2rong2}
  \antonymref{缓缓}{huan3huan3}
  \antonymref{缓慢}{huan3man4}
  \antonymref{慢慢}{man4man4}
  \antonymref{徐徐}{xu2xu2}
  \end{Phonetics}
\end{Entry}

\begin{Entry}{赶脚}{10,11}{⾛,⾁}
  \begin{Phonetics}{赶脚}{gan3jiao3}
    \definition{v.}{contratar para transportar mercadorias com seu próprio burro, mula, etc. | transportar mercadorias para ganhar a vida (especialmente usando um burro) | trabalhar como carreteiro ou carregador}
  \end{Phonetics}
\end{Entry}

\begin{Entry}{赶跑}{10,12}{⾛,⾜}
  \begin{Phonetics}{赶跑}{gan3pao3}
    \definition{v.}{ir embora | forçar a saída | repelir}
  \synonymref{追赶}{zhui1gan3}
  \end{Phonetics}
\end{Entry}

\begin{Entry}{赶集}{10,12}{⾛,⾫}
  \begin{Phonetics}{赶集}{gan3ji2}
    \definition{v.}{ir a uma feira; ir a uma feira rural/mercado local; ir ao mercado para comprar ou vender coisas ou para passear}
  \end{Phonetics}
\end{Entry}

\begin{Entry}{赶路}{10,13}{⾛,⾜}
  \begin{Phonetics}{赶路}{gan3lu4}
    \definition{v.}{apressar"-se em sua jornada; acelerar a viagem para chegar mais cedo}
  \synonymref{打量}{da3liang5}
  \synonymref{走路}{zou3lu4}
  \antonymref{住宿}{zhu4su4}
  \end{Phonetics}
\end{Entry}

%%%%%%%%%% 起 %%%%%%%%%%
\subsection*{起}\addcontentsline{loh}{figure}{起}

\begin{Entry}{起}{10}{⾛}
  \begin{Phonetics}{起}{qi3}[][HSK 1]
    \definition{clas.}{caso; instância | lote; grupo}
    \definition{prep.}{de; colocado antes de uma palavra de tempo ou lugar, indica um ponto de partida | por; colocado antes de uma palavra de lugar, indica um lugar por onde passou}
    \definition{v.}{levantar"-se; ficar de pé| iniciar; lançar; deicar a posição original | subir; ascender | aparecer; levantar; crescer (bolhas, protuberâncias, brotoeja) | puxar para cima; puxar para fora; tirar o que está guardado ou incorporado | crescer; aumentar | esboçar; elaborar | construir; montar; estabelecer | receber (comprovante) | começar; iniciar; combina com 从 e 由; indica quando, onde e quem começou | buscar; pegar; usado após um verbo, indica movimento para cima | indicar se alguém tem força suficiente ou não; usado após um verbo, indica que a força é suficiente ou insuficiente | indicar que a ação envolve alguém ou algo; equivalente a 及 ou 到 | começar; iniciar; usado depois de um verbo, indica o início de uma ação | juntar; implodir; (informal) usado depois de um verbo, para unir coisas ou fechá"-las}
  \seealsoref{从}{cong2}
  \seealsoref{到}{dao4}
  \seealsoref{及}{ji2}
  \seealsoref{由}{you2}
  \synonymref{崇}{chong2}
  \synonymref{次}{ci4}
  \synonymref{件}{jian4}
  \synonymref{建}{jian4}
  \synonymref{拟}{ni3}
  \synonymref{群}{qun2}
  \synonymref{设}{she4}
  \synonymref{升}{sheng1}
  \synonymref{始}{shi3}
  \antonymref{闭}{bi4}
  \antonymref{落}{la4}
  \antonymref{落}{lao4}
  \antonymref{了}{liao3}
  \antonymref{落}{luo4}
  \antonymref{收}{shou1}
  \antonymref{止}{zhi3}
  \end{Phonetics}
\end{Entry}

\begin{Entry}{起飞}{10,3}{⾛,⾶}
  \begin{Phonetics}{起飞}{qi3fei1}[][HSK 2]
    \definition{v.}{decolar; levantar voo | crescer rapidamente; decolar; disparar; metáfora para o rápido desenvolvimento de negócios, economia, etc.}
  \synonymref{升高}{sheng1gao1}
  \synonymref{升起}{sheng1qi3}
  \antonymref{降落}{jiang4luo4}
  \antonymref{下降}{xia4jiang4}
  \end{Phonetics}
\end{Entry}

\begin{Entry}{起伏}{10,6}{⾛,⼈}
  \begin{Phonetics}{起伏}{qi3fu2}[][HSK 7-9]
    \definition{v.}{subir e descer; ondular;  subida e descida contínuas | (emoções, relações, etc.) flutuar; oscilar; aumentar e diminuir; essa metáfora descreve a oscilação de emoções, relacionamentos, etc.}
  \end{Phonetics}
\end{Entry}

\begin{Entry}{起初}{10,7}{⾛,⾐}
  \begin{Phonetics}{起初}{qi3chu1}[][HSK 7-9]
    \definition{adv.}{princípio; início; inicial (frequentemente contrastado com 后来 ou 现在)}
  \seealsoref{后来}{hou4lai2}
  \seealsoref{现在}{xian4zai4}
  \end{Phonetics}
\end{Entry}

\begin{Entry}{起劲}{10,7}{⾛,⼒}
  \begin{Phonetics}{起劲}{qi3jin4}[][HSK 7-9]
    \definition{adj.}{vigoroso; animado; entusiasmado; enérgico; descreve um alto nível de entusiasmo por fazer algo}
  \seealsoref{起劲儿}{qi3jin4r5}
  \end{Phonetics}
\end{Entry}

\begin{Entry}{起劲儿}{10,7,2}{⾛,⼒,⼉}
  \begin{Phonetics}{起劲儿}{qi3jin4r5}
    \definition{v.}{ser enérgico}
  \end{Phonetics}
\end{Entry}

\begin{Entry}{起床}{10,7}{⾛,⼴}
  \begin{Phonetics}{起床}{qi3/chuang2}[][HSK 1]
    \definition{v.+compl.}{levantar"-se; sair da cama; acordar e sair da cama (geralmente pela manhã); levantar"-se da posição sentada, deitada ou deitada de bruços, ou sentar"-se a partir da posição deitada}
  \synonymref{起来}{qi3/lai2}
  \antonymref{睡觉}{shui4/jiao4}
  \end{Phonetics}
\end{Entry}

\begin{Entry}{起来}{10,7}{⾛,⽊}
  \begin{Phonetics}{起来}{qi3/lai2}[][HSK 1]
    \definition{v.+compl.}{levantar"-se; passar de posições como deitado, sentado ou ajoelhado para ficar em pé | levantar"-se; sair da cama | levantar"-se; revoltar"-se; rebelar"-se; refere"-se a ascensão, surgimento, levantamento, etc.}
  \seealsoref{上来}{shang4lai5}
  \synonymref{上来}{shang4lai5}
  \end{Phonetics}
  \begin{Phonetics}{起来}{qi3lai5}
    \definition{v.aux.}{usado depois de um verbo para indicar movimento ascendente}
  \end{Phonetics}
  \begin{Phonetics}{起来}{qi5lai2}
    \definition{v.}{descrever resultados, retratar comportamentos, transmitir movimento}
  \end{Phonetics}
\end{Entry}

\begin{Entry}{起步}{10,7}{⾛,⽌}
  \begin{Phonetics}{起步}{qi3bu4}[][HSK 7-9]
    \definition{v.}{começar; mover"-se; começar a dar os primeiros passos | começar; dar início; metaforicamente, significa o início de (uma carreira, emprego, etc.)}
  \end{Phonetics}
\end{Entry}

\begin{Entry}{起诉}{10,7}{⾛,⾔}
  \begin{Phonetics}{起诉}{qi3su4}[][HSK 6]
    \definition{v.}{processar; entrar com uma ação judicial}
  \end{Phonetics}
\end{Entry}

\begin{Entry}{起到}{10,8}{⾛,⼑}
  \begin{Phonetics}{起到}{qi3dao4}[][HSK 5]
    \definition{v.}{ter (um efeito motivador, etc.); desempenhar (um papel estabilizador, etc.)}
  \end{Phonetics}
\end{Entry}

\begin{Entry}{起码}{10,8}{⾛,⽯}
  \begin{Phonetics}{起码}{qi3ma3}[][HSK 5]
    \definition{adj.}{mínimo; elementar; rudimentar}
    \definition{adv.}{minimamente; pelo menos}
  \end{Phonetics}
\end{Entry}

\begin{Entry}{起点}{10,9}{⾛,⽕}
  \begin{Phonetics}{起点}{qi3dian3}[][HSK 6]
    \definition[个]{s.}{ponto de partida (para o tempo ou local do início de algo); o lugar ou hora de início | ponto de partida (para o nível ou base de algo feito inicialmente); refere"-se especificamente ao ponto de partida designado em um evento de pista}
  \end{Phonetics}
\end{Entry}

\begin{Entry}{起草}{10,9}{⾛,⾋}
  \begin{Phonetics}{起草}{qi3/cao3}[][HSK 7-9]
    \definition{v.+compl.}{elaborar; redigir; preparar um rascunho}
  \end{Phonetics}
\end{Entry}

\begin{Entry}{起程}{10,12}{⾛,⽲}
  \begin{Phonetics}{起程}{qi3cheng2}[][HSK 7-9]
    \definition{v.}{partir; sair; iniciar uma viagem}
  \seealsoref{启程}{qi3cheng2}
  \end{Phonetics}
\end{Entry}

\begin{Entry}{起跑线}{10,12,8}{⾛,⾜,⽷}
  \begin{Phonetics}{起跑线}{qi3pao3xian4}[][HSK 7-9]
    \definition{s.}{linha de partida (em uma corrida de revezamento); a linha de largada (de uma corrida); a linha de partida em provas de atletismo; metaforicamente, o ponto de partida para o trabalho, estudo, etc.}
  \end{Phonetics}
\end{Entry}

\begin{Entry}{起源}{10,13}{⾛,⽔}
  \begin{Phonetics}{起源}{qi3yuan2}[][HSK 7-9]
    \definition{s.}{a origem; a origem mais remota das coisas}
    \definition{v.}{originar"-se; ter origem em; começar a acontecer; começar a aparecer}
  \end{Phonetics}
\end{Entry}

\begin{Entry}{起跳}{10,13}{⾛,⾜}
  \begin{Phonetics}{起跳}{qi3tiao4}
    \definition{v.}{Atletismo: decolar (no início de um salto) | (preço, salário, etc.) começar (de um determinado nível)}
  \end{Phonetics}
\end{Entry}

%%%%%%%%%% 趁 %%%%%%%%%%
\subsection*{趁}\addcontentsline{loh}{figure}{趁}

\begin{Entry}{趁}{12}{⾛}
  \begin{Phonetics}{趁}{chen4}[][HSK 7-9]
    \definition{prep.}{aproveitar"-se de; tirar vantagem de (tempo, oportunidade, etc.); indica o tempo e as condições de uso}
    \definition{v.}{ser rico em; possuir}
  \end{Phonetics}
\end{Entry}

\begin{Entry}{趁早}{12,6}{⾛,⽇}
  \begin{Phonetics}{趁早}{chen4zao3}[][HSK 7-9]
    \definition{adv.}{o mais cedo possível; antes que seja tarde demais; na primeira oportunidade}
  \end{Phonetics}
\end{Entry}

\begin{Entry}{趁机}{12,6}{⾛,⽊}
  \begin{Phonetics}{趁机}{chen4ji1}[][HSK 7-9]
    \definition{adv./adv.}{aproveitar a ocasião; aproveitar a oportunidade}
  \end{Phonetics}
\end{Entry}

\begin{Entry}{趁着}{12,11}{⾛,⽬}
  \begin{Phonetics}{趁着}{chen4zhe5}[][HSK 7-9]
    \definition{v.}{aproveitar"-se de; tirar vantagem de}
  \end{Phonetics}
\end{Entry}

%%%%%%%%%% 超 %%%%%%%%%%
\subsection*{超}\addcontentsline{loh}{figure}{超}

\begin{Entry}{超}{12}{⾛}
  \begin{Phonetics}{超}{chao1}[][HSK 6]
    \definition{adj.}{super; extremamente; maior (ou menor) que o padrão geral}
    \definition{v.}{exceder; ultrapassar; vir para a frente por trás; prevalecer | transcender; ir além; não ser sujeito a certas restrições; ir além de um certo intervalo | exceder; superar; exceder o limite prescrito}
  \end{Phonetics}
\end{Entry}

\begin{Entry}{超车}{12,4}{⾛,⾞}
  \begin{Phonetics}{超车}{chao1/che1}[][HSK 7-9]
    \definition{v.+compl.}{ultrapassar (um veículo); passar por um veículo que trafega na mesma direção}
  \end{Phonetics}
\end{Entry}

\begin{Entry}{超出}{12,5}{⾛,⼐}
  \begin{Phonetics}{超出}{chao1chu1}[][HSK 6]
    \definition{v.}{exceder; ultrapassar; ir além (de uma certa quantidade ou intervalo)}
  \end{Phonetics}
\end{Entry}

\begin{Entry}{超市}{12,5}{⾛,⼱}
  \begin{Phonetics}{超市}{chao1shi4}[][HSK 2]
    \definition[家]{s.}{supermercado; abreviação de 超级市场}
  \seealsoref{超级市场}{chao1ji2shi4chang3}
  \synonymref{商场}{shang1chang3}
  \synonymref{商店}{shang1dian4}
  \end{Phonetics}
\end{Entry}

\begin{Entry}{超级}{12,6}{⾛,⽷}
  \begin{Phonetics}{超级}{chao1ji2}[][HSK 3]
    \definition{adj.}{super; além do nível geral}
    \definition{pref.}{super-; ultra-; hiper-}
  \synonymref{出众}{chu1zhong4}
  \synonymref{顶级}{ding3ji2}
  \synonymref{顶尖}{ding3jian1}
  \synonymref{一流}{yi4liu2}
  \antonymref{平常}{ping2chang2}
  \antonymref{普通}{pu3tong1}
  \antonymref{一般}{yi4ban1}
  \end{Phonetics}
\end{Entry}

\begin{Entry}{超级市场}{12,6,5,6}{⾛,⽷,⼱,⼟}
  \begin{Phonetics}{超级市场}{chao1ji2shi4chang3}
    \definition[个,间,所,家]{s.}{supermercado; hipermercado}
  \end{Phonetics}
\end{Entry}

\begin{Entry}{超过}{12,6}{⾛,⾡}
  \begin{Phonetics}{超过}{chao1/guo4}[][HSK 2]
    \definition{v.+compl.}{ultrapassar; superar (algo ou alguém); passar de trás para a frente de alguém ou algo | exceder; ser mais do que; ultrapassar (um padrão)}
  \seealsoref{超越}{chao1yue4}
  \synonymref{超出}{chao1chu1}
  \synonymref{超越}{chao1yue4}
  \synonymref{赶上}{gan3shang4}
  \synonymref{领先}{ling3/xian1}
  \synonymref{越过}{yue4/guo4}
  \antonymref{不及}{bu4ji2}
  \antonymref{不如}{bu4ru2}
  \antonymref{落后}{luo4/hou4}
  \antonymref{滞后}{zhi4hou4}
  \end{Phonetics}
\end{Entry}

\begin{Entry}{超声}{12,7}{⾛,⼠}
  \begin{Phonetics}{超声}{chao1sheng1}
    \definition{adj.}{ultrasônico; supersônico}
    \definition{s.}{ultrassom; supersom}
  \end{Phonetics}
\end{Entry}

\begin{Entry}{超前}{12,9}{⾛,⼑}
  \begin{Phonetics}{超前}{chao1qian2}[][HSK 7-9]
    \definition{v.}{transcender os tempos; estar à frente dos tempos; estar à frente do seu tempo | superar os antecessores; liderar; assumir a liderança}
  \end{Phonetics}
\end{Entry}

\begin{Entry}{超标}{12,9}{⾛,⽊}
  \begin{Phonetics}{超标}{chao1/biao1}[][HSK 7-9]
    \definition{v.+compl.}{exceder uma cota (ou padrão); exceder o limite prescrito}
  \end{Phonetics}
\end{Entry}

\begin{Entry}{超速}{12,10}{⾛,⾡}
  \begin{Phonetics}{超速}{chao1su4}[][HSK 7-9]
    \definition{s.}{Física: hipervelocidade | excesso de velocidade}
    \definition{v.}{exceder o limite de velocidade}
  \end{Phonetics}
\end{Entry}

\begin{Entry}{超越}{12,12}{⾛,⾛}
  \begin{Phonetics}{超越}{chao1yue4}[][HSK 5]
    \definition{v.}{ultrapassar; superar; passar por cima; transcender}
  \end{Phonetics}
\end{Entry}

%%%%%%%%%% 越 %%%%%%%%%%
\subsection*{越}\addcontentsline{loh}{figure}{越}

\begin{Entry}{越}{12}{⾛}
  \begin{Phonetics}{越}{yue4}[][HSK 2]
    \definition{adj.}{superior; excede ou ultrapassa o ordinário}
    \definition{adv.}{quanto mais\dots mais; usados juntos, eles formam o formato de ``越…越…'' para indicar que o grau de uma situação se torna mais sério à medida que se desenvolve; ``成年…'' para indicar que o grau de uma situação se torna mais sério à medida que o tempo passa}
    \definition{v.}{passar por cima; pular; cruzar | exceder; ultrapassar | estar em um tom alto; estar animado | saquear; pilhar; expoliar; apreender; roubar | passar; passar através; atravessar}
  \seealsoref{越……越……}{yue4 yue4}
  \seealsoref{越来越……}{yue4lai2yue4}
  \synonymref{昂}{ang2}
  \synonymref{昂}{ang2}
  \synonymref{超}{chao1}
  \synonymref{超}{chao1}
  \synonymref{夺}{duo2}
  \synonymref{夺}{duo2}
  \synonymref{翻}{fan1}
  \synonymref{翻}{fan1}
  \synonymref{劫}{jie2}
  \synonymref{劫}{jie2}
  \synonymref{经}{jing1}
  \synonymref{经}{jing1}
  \synonymref{跨}{kua4}
  \synonymref{跨}{kua4}
  \synonymref{抢}{qiang3}
  \synonymref{抢}{qiang3}
  \synonymref{跳}{tiao4}
  \synonymref{跳}{tiao4}
  \synonymref{扬}{yang2}
  \synonymref{扬}{yang2}
  \end{Phonetics}
\end{Entry}

\begin{Entry}{越发}{12,5}{⾛,⼜}
  \begin{Phonetics}{越发}{yue4fa1}[][HSK 7-9]
    \definition{adv.}{ainda mais; tanto mais; cada vez mais; isso indica um grau adicional, equivalente a 更加 | quanto mais\dots quantomais\dots; ecoa o 越 ou 或 anterior e sua função é a mesma de 越……越…… (usado quando duas ou mais orações se repetem)}
  \seealsoref{更加}{geng4jia1}
  \seealsoref{或}{huo4}
  \seealsoref{越}{yue4}
  \seealsoref{越……越……}{yue4 yue4}
  \synonymref{更加}{geng4jia1}
  \synonymref{尤其}{you2qi2}
  \antonymref{依旧}{yi1jiu4}
  \end{Phonetics}
\end{Entry}

\begin{Entry}{越过}{12,6}{⾛,⾡}
  \begin{Phonetics}{越过}{yue4/guo4}[][HSK 7-9]
    \definition{v.+compl.}{cruzar; superar; negociar; transpor obstáculos, fronteiras, etc. | exceder; ultrapassar}
  \synonymref{超过}{chao1/guo4}
  \synonymref{超出}{chao1chu1}
  \synonymref{超越}{chao1yue4}
  \synonymref{穿过}{chuan1guo4}
  \synonymref{跨越}{kua4yue4}
  \synonymref{突出}{tu1/chu1}
  \end{Phonetics}
\end{Entry}

\begin{Entry}{越来越……}{12,7,12}{⾛,⽊,⾛}
  \begin{Phonetics}{越来越……}{yue4lai2yue4}[][HSK 2]
    \definition{adv.}{cada vez mais\dots; isso significa que o grau de algo se aprofunda à medida que o tempo passa}
  \seealsoref{越……越……}{yue4 yue4}
  \synonymref{愈来愈}{yu4lai2yu4}
  \synonymref{越……越……}{yue4 yue4}
  \end{Phonetics}
\end{Entry}

\begin{Entry}{越……越……}{12,12}{⾛,⾛}
  \begin{Phonetics}{越……越……}{yue4 yue4}
    \definition{expr.}{quanto mais\dots tanto mais\dots}
  \end{Phonetics}
\end{Entry}

\begin{Entry}{越障}{12,13}{⾛,⾩}
  \begin{Phonetics}{越障}{yue4zhang4}
    \definition{s.}{corrida de obstáculos; corrida com obstáculos | pista de obstáculos para treinamento de tropas}
    \definition{v.}{superar obstáculos}
  \end{Phonetics}
\end{Entry}

\begin{Entry}{越境}{12,14}{⾛,⼟}
  \begin{Phonetics}{越境}{yue4jing4}
    \definition{v.}{cruzar a fronteira ilegalmente; entrar ou sair clandestinamente de um país}
  \end{Phonetics}
\end{Entry}

%%%%%%%%%% 趋 %%%%%%%%%%
\subsection*{趋}\addcontentsline{loh}{figure}{趋}

\begin{Entry}{趋}{12}{⾛}
  \begin{Phonetics}{趋}{qu1}
    \definition{v.}{apressar"-se | tender para; tender a se tornar | (ganso, cobra, etc.) estalar a cabeça e morder as pessoas}
  \end{Phonetics}
\end{Entry}

\begin{Entry}{趋于}{12,3}{⾛,⼆}
  \begin{Phonetics}{趋于}{qu1yu2}[][HSK 7-9]
    \definition{v.}{tender a}
  \end{Phonetics}
\end{Entry}

\begin{Entry}{趋势}{12,8}{⾛,⼒}
  \begin{Phonetics}{趋势}{qu1shi4}[][HSK 4]
    \definition{s.}{rumo; tendência; direção; impulso das coisas que se movem em uma direção ou outra}
  \end{Phonetics}
\end{Entry}

%%%%%%%%%% 趟 %%%%%%%%%%
\subsection*{趟}\addcontentsline{loh}{figure}{趟}

\begin{Entry}{趟}{15}{⾛}
  \begin{Phonetics}{趟}{tang1}
    \definition{v.}{atravessar; andar na grama ou onde não haja caminho | usar arados, capinadores, etc. para virar o solo e remover ervas daninhas | vadear; atravessar a vau; caminhar por águas rasas}[我们趟水去那小岛。===Nós vadeamos até a ilha.]
  \end{Phonetics}
  \begin{Phonetics}{趟}{tang4}[][HSK 6]
    \definition{clas.}{usado para o número de vezes de viagens de ida e volta |  usado para coisas dispostas em fileiras ou tiras | usado para a programação de veículos, navios, etc. que circulam em uma determinada ordem | usado em conjuntos de movimentos de artes marciais}
    \definition{s.}{marcha; procissão; jornada; viagem}
  \end{Phonetics}
\end{Entry}

%%%%%%%%%% 趣 %%%%%%%%%%
\subsection*{趣}\addcontentsline{loh}{figure}{趣}

\begin{Entry}{趣}{15}{⾛}
  \begin{Phonetics}{趣}{qu4}
    \definition{adj.}{interessante}
    \definition{s.}{interesse; deleite; diversão; passatempo | propósito; inclinação}
  \end{Phonetics}
\end{Entry}

\begin{Entry}{趣味}{15,8}{⾛,⼝}
  \begin{Phonetics}{趣味}{qu4wei4}[][HSK 7-9]
    \definition{adj.}{agradável; interessante; atraente; simpático}
    \definition[种]{s.}{interesse; deleite; passatempo}
  \end{Phonetics}
\end{Entry}

%%%%% EOF %%%%%

